\hypertarget{4-results}{%
\section{Results}\label{4-results}}

\hypertarget{41-des-validation}{%
\subsection{DES validation}\label{41-des-validation}}

The reconstructed DES was first validated against the deterministic
baseline (Cf0, single shift S=1) reported in Garrido-Rios (2017, Table
6.10)~\cite{garrido2017}. Using the official thesis basis of 8,064 hours per year (48
operational weeks), the simulation delivers 733,621 rations per year on
average, compared with the thesis reference of 767,592 rations, yielding
a relative gap of $-4.43\%$. This is comfortably within the accepted
$\pm 15\%$ validation threshold and supports the structural fidelity of
the rebuilt DES. A secondary Gregorian-basis check (8,760 hours per
year) produces 796,940 rations per year ($+3.82\%$ relative gap),
confirming that the discrepancy is driven by year-basis conventions
rather than model error.

Under deterministic conditions with risks disabled, the simulation
produces 734,458 rations per year and delivers 733,621, with a fill rate
of 99.3\% and only 41 backorders accumulated over the 20-year horizon.
When stochastic risks are enabled at the current level, average annual
delivery falls to 677,750 rations ($-7.6\%$ relative to the
deterministic baseline), the fill rate drops to 68.3\%, and total
backorders rise to 1,825. Under the increased risk configuration used
for the benchmark, average annual delivery decreases further to 549,250
rations, fill rate falls to 45.6\%, and total backorders increase to
3,132. The DES therefore shows a clear transition from a well-calibrated
nominal system to a fragile network once disruption intensity escalates.

\hypertarget{42-track-a}{%
\subsection{Track A: thesis-grounded benchmark (negative
result)}\label{42-track-a}}

Under the thesis-grounded Track A contract, we first ask whether a
regime-conditioned oracle over the static decision space (inventory
buffer fraction and shift level) exposes any exploitable structure. A
dense grid search over five risk families, five frequency multipliers,
and three impact multipliers finds a real but small oracle headroom: the
best regime-varying combination of static actions beats the single best
constant policy by +0.000176 to +0.000296 Excel ReT (95\% CI lower bound
positive in the full grid), with the optimal action changing across
regimes. This headroom is small but not zero, so the relevant question is
whether a learned policy can convert it.

No tested learner does. Continuous PPO, PPO warm-started via behavior
cloning on the oracle action table, and PPO trained under a family-conflict
action contract (\texttt{per\_op\_conflict}, reward
\texttt{ReT\_excel\_plus\_cvar}, 5 seeds) all fail to close the gap under
paired common-random-number evaluation against the same dense static
grid. The closest approach --- behavior-cloning-warm-started PPO under
the family-conflict contract --- misses the best static comparator by
$-6.99\times10^{-6}$ Excel ReT (0.155247 vs.\ 0.155254), a difference two
orders of magnitude smaller than the measured oracle headroom, and is
Pareto-dominated by 10 static configurations on the joint Excel
ReT--resource frontier. We therefore report Track A as a boundary
characterization: under dense common-random-number static frontiers and
the tested reward/algorithm combinations, no publishable dynamic-control
signal survives, even though a small regime-conditioned oracle headroom
is measurable. This is not claimed as a proof of mathematical
impossibility, only as the outcome of the tested learners under fair,
paired evaluation.

The structural explanation is the downstream distribution bottleneck:
Op10--Op12 transport throughput ($\approx$2,500 rations/day)
approximately equals demand, so extra upstream production capacity from
higher shift levels yields sharply diminishing returns once a modest
buffer is in place. Because Track A does not expose downstream dispatch
control to the agent, the active constraint lies outside the
controllable interface --- motivating the Track B extension in
Section~4.3.

\hypertarget{43-track-b}{%
\subsection{Track B: extended action space (positive
result)}\label{43-track-b}}

Track B extends the action contract (\texttt{track\_b\_v1}, 8 dimensions)
by adding downstream dispatch multipliers at Op10 and Op12 on top of the
Track A buffer/shift variables, while preserving the DES backbone and
reward family (\texttt{control\_v1}). Table~\ref{tab:track_b_main} presents
the canonical confirmatory result (5 seeds $\times$ 60{,}000 timesteps
$\times$ 12 evaluation episodes, \texttt{adaptive\_benchmark\_v2} risk
level, $h{=}104$ weeks), comparing the learned policy against the best of a
dense grid of 147 static dispatch policies (3 shift levels $\times$ 7 Op10
multipliers $\times$ 7 Op12 multipliers) under paired common-random-number
(CRN) evaluation.

\begin{table}[htbp]
\centering
\caption{Track B main result: PPO vs.\ best dense static comparator (\texttt{S2}, Op10$\times$2.00, Op12$\times$1.50), paired CRN evaluation (5 training seeds $\times$ 12 CRN episodes). Episode-level pooled CI95 in brackets, reported descriptively; primary inference is seed-clustered (all five per-seed mean deltas positive; seed-clustered bootstrap CI95 for the Excel ReT delta $[+0.000414, +0.000438]$).}
\label{tab:track_b_main}
\small
\begin{tabular}{lccc}
\toprule
Metric & PPO & Best static & $\Delta$ [CI95] \\
\midrule
Excel ReT & 0.005893 & 0.005466 & +0.000426 [0.000389, 0.000463] \\
Order-level ReT & 0.005666 & 0.005251 & +0.000415 [0.000378, 0.000450] \\
Flow fill rate & 0.961 & 0.669 & +0.293 [0.281, 0.304] \\
Rolling fill (4w) & 0.996 & 0.909 & +0.087 [0.034, 0.154] \\
CTj p99 (hours) & 1{,}207 & 8{,}113 & $-6{,}906$ [$-7{,}276$, $-6{,}538$] \\
Service-loss AUC/order & 113{,}170 & 1{,}125{,}480 & $-1{,}012{,}310$ \\
Assembly cost index & 0.682 & 0.667 & +0.015 [$-0.048$, 0.015] (n.s.) \\
\bottomrule
\end{tabular}
\end{table}

PPO improves Garrido/Excel ReT by 0.000426 over the best of 147 dense
static comparators. Because the 60 CRN pairs share only 16 distinct
evaluation scenarios across seeds, we treat the training seed as the
inferential unit: all five per-seed mean deltas are individually positive
(range $+0.000408$ to $+0.000447$), and the seed-clustered bootstrap CI
excludes zero. The gain is accompanied by matching, larger-magnitude
improvements in flow fill rate, rolling fill, and
recovery-tail cycle time (CTj p99 falls from 8{,}113 to 1{,}207 hours).
The policy is Pareto non-dominated on the ReT--cost--tail--flow frontier:
no static policy in the dense grid jointly matches or exceeds PPO on Excel
ReT, cost, CTj p99, and flow fill simultaneously. It is \emph{not},
however, a strict cost win: the assembly cost index is marginally higher
for PPO than the best static comparator by ReT, and this difference is not
statistically distinguishable from zero (CI95 crosses zero). We therefore
report the result as non-dominated, not resource-efficient.

\begin{figure}[htbp]
\centering
\includegraphics[width=0.78\linewidth]{fig4_pareto_ret_tail_ctj.png}
\caption{Track B ReT--tail Pareto view from the final Q1 bundle. PPO is not a strict assembly-cost winner, but it is non-dominated on the resilience/tail frontier relative to the dense static dispatch family.}
\label{fig:pareto_ret_tail}
\end{figure}

A companion masking experiment addresses whether this advantage could
instead reflect the learned policy exploiting privileged simulator
information --- specifically, a one-hot encoding of the true disruption
regime and forecast features derived from the true regime-transition
matrix, both present in the observation vector but not available to any
static comparator. Retraining PPO with these seven fields masked to zero
(Section~\ref{45-privileged-obs}) retains approximately 95\% of the
canonical advantage (+0.000395 to +0.000403 across two independent
training runs on separate infrastructure), and a non-learning policy given
direct access to the true regime signal (Section~\ref{45-privileged-obs})
still falls far short of PPO. The advantage is not an artifact of
privileged observation access.

\hypertarget{44-causal-ablation}{%
\subsection{Causal ablation: action-space decomposition}\label{44-causal-ablation}}

To test whether the Track B advantage reflects access to the downstream
bottleneck specifically, or simply a larger action space, we run a
5-seed $\times$ 60{,}000-timestep ablation over the 8-dimensional
\texttt{track\_b\_v1} contract, holding reward (\texttt{control\_v1}),
observation version (v7), and risk level (\texttt{adaptive\_benchmark\_v2})
fixed. Three configurations are compared: \emph{joint} (the full 8D
contract), \emph{downstream\_only} (shift frozen at S2, only Op10/Op12
dispatch controllable), and \emph{shift\_only} (dispatch frozen at neutral,
only shift controllable).

\begin{table}[htbp]
\centering
\caption{Action-space ablation, 5 seeds $\times$ 60k timesteps. $\Delta$ is order-level ReT versus each arm's own best evaluated comparator in that arm, including heuristic comparators.}
\label{tab:ablation}
\small
\begin{tabular}{lccc}
\toprule
Action config & PPO ReT & Best in-arm comparator & $\Delta$ \\
\midrule
Joint (8D) & 0.005587 & 0.005221 & +0.000367 \\
\textbf{Downstream-only} & \textbf{0.005678} & 0.005249 & \textbf{+0.000429} \\
Shift-only & 0.005607 & 0.005231 & +0.000377 \\
\bottomrule
\end{tabular}
\end{table}

The gain does not scale with the number of controllable dimensions: a
downstream-only policy captures a \emph{larger} advantage over its own best
evaluated comparator (+0.000429) than the full joint 8-dimensional contract
(+0.000367), with shift-only capturing a comparable advantage
(+0.000377) to the joint contract. This rules out the confound that Track
B wins simply because it gives the learner more actuators; the effect
concentrates specifically in access to downstream dispatch. We retain the
joint 8D contract as the paper's canonical result (Section~4.3, trained at
full 60k$\times$5-seed scale with matched hyperparameters to Track A) and
treat the ablation as mechanism evidence rather than a recommendation to
drop the shift dimensions.\footnote{The ablation runs used learning rate
$1\times10^{-4}$, differing from the canonical run's $3\times10^{-4}$;
absolute ReT values are therefore not directly comparable to
Table~\ref{tab:track_b_main}, but the within-ablation comparison across
arms (identical hyperparameters in all three) is valid.}

\begin{figure}[htbp]
\centering
\includegraphics[width=0.78\linewidth]{fig6_action_space_ablation}
\caption{E4 action-space ablation. Deltas are computed against each arm's own best evaluated comparator, including heuristics, and support the mechanism claim that downstream dispatch authority is central to the Track B gain.}
\label{fig:ablation}
\end{figure}

\hypertarget{45-privileged-obs}{%
\subsection{Ruling out privileged observation access}\label{45-privileged-obs}}

Observation version v7 includes, alongside operationally observable state
(backlog, queue pressure, rolling fill history), a one-hot encoding of the
simulator's true disruption regime and two forecast features derived from
the true regime-transition matrix. Because static and heuristic
comparators cannot condition on these fields, a skeptical reading of
Section~4.3 is that PPO's advantage reflects exploiting privileged
simulator state rather than learning genuine closed-loop control. We test
this from both directions.

\emph{Does PPO need the privileged fields?} We retrain PPO with the
regime one-hot and forecast fields masked to zero, holding all other
hyperparameters at their canonical values, and reproduce the run
independently on two separate compute environments. The masked policy
retains the great majority of the canonical advantage. Against the static
comparator family, the masked policy's order-level ReT advantage is
$+0.000403$ on the OVH replication and approximately $+0.000395$ in the
Kaggle replication, versus $+0.000415$ for the canonical, unmasked policy.
Including the strongest heuristic comparator narrows but does not remove
the margin; the OVH masked PPO mean is 0.005613 versus 0.005245 for the
best heuristic.

\emph{Does a non-learning policy with the same privileged access match
PPO?} We construct a regime-conditioned static lookup table --- the best
constant dispatch policy per observed disruption regime, fit directly on
the true regime signal PPO could in principle exploit --- and evaluate it,
six reactive dispatch heuristics, and 75 constant-dispatch policies,
CRN-paired against the canonical frozen PPO checkpoint. The fitted regime
table improves on the single best static constant by only $+0.0000277$,
a fraction of PPO's $+0.000427$ advantage over the same comparator; PPO
beats all 75 static candidates, all six heuristics, and the fitted regime
table individually, with 95\% confidence intervals excluding zero in
every one of 60 comparisons.

\begin{figure}[htbp]
\centering
\includegraphics[width=0.78\linewidth]{fig3_gap_decomposition}
\caption{E1 gap decomposition. Direct access to the true disruption regime improves a non-learning static table only modestly, while PPO remains far above the common static, heuristic, and regime-table comparators.}
\label{fig:gap_decomposition}
\end{figure}

Together these results indicate the Track B advantage is not explained by
access to privileged simulator state: PPO does not require it to win, and
a non-learning policy that has it still cannot approach PPO's
performance.

\begin{table}[htbp]
\centering
\caption{Privileged-observation defense. E1 values use the canonical CRN comparison; the masked PPO row is from the E2 retrain with regime and forecast fields zeroed.}
\label{tab:privileged_defense}
\small
\begin{tabular}{p{0.30\linewidth}c p{0.40\linewidth}}
\toprule
Policy family & ReT & Interpretation \\
\midrule
Common dense static & 0.005466 & Best 147-cell constant dispatch comparator \\
Best E1 heuristic & 0.005436 & Reactive rule family does not match PPO \\
Regime-conditioned table & 0.005494 & True-regime lookup gives only a small gain \\
Masked PPO & 0.005613 & Wins after seven privileged fields are removed \\
Canonical PPO & 0.005893 & Full Track B policy \\
\bottomrule
\end{tabular}
\end{table}

\hypertarget{46-reward-sweep}{%
\subsection{Reward sensitivity under Track B}\label{46-reward-sweep}}

To assess whether the Track B advantage depends on the specific choice of
training reward, we screen 18 reward/observation combinations under the
current 8-dimensional action contract: reward modes \texttt{control\_v1},
\texttt{ReT\_excel\_plus\_cvar} (tail weight $\alpha \in \{0.05, 0.1,
0.2\}$), \texttt{ReT\_tail\_v2}, and \texttt{ReT\_garrido2024\_train},
crossed with observation versions v7, v8, and v9, each trained for
2 seeds $\times$ 40{,}000 timesteps (a screening budget, smaller than the
5-seed $\times$ 60{,}000-timestep confirmatory scale used for the
canonical result) and evaluated CRN-paired against the best dense static
comparator.

\textbf{Every one of the 18 cells shows a positive Excel ReT delta}
(range $+0.000195$ to $+0.000452$), and every cell registers a directional
win on the primary metric. The advantage does not depend on the specific
reward formulation, the CVaR tail weight, or the observation version
tested. This screen is not a substitute for a confirmatory-scale sweep ---
promotion decisions in this screen selected 12 of 18 cells for further
scaling, using a threshold on Excel ReT and CVaR05 improvement --- but the
uniform sign and consistent magnitude across 18 independently-trained
configurations is evidence against the concern that the canonical result
depends on a fortuitously tuned reward.

\hypertarget{47-cross-scenario}{%
\subsection{Cross-scenario and cross-horizon generalization}\label{47-cross-scenario}}

To test whether the Track B advantage is specific to the single
\texttt{adaptive\_benchmark\_v2} stress cell used for the canonical result,
we evaluate the frozen, canonical PPO checkpoints (no retraining) against a
fixed set of static comparators, CRN-paired, across three
Garrido-native risk levels (\texttt{current}, \texttt{increased},
\texttt{severe}) and two horizons ($h{=}52$ and $h{=}104$ weeks; 5 seeds
$\times$ 12 episodes per cell).

\begin{table}[htbp]
\centering
\caption{Cross-regime/horizon generalization matrix (frozen PPO vs. best comparator, CRN-paired).}
\label{tab:cross_scenario}
\begin{tabular}{llrrr}
\toprule
Horizon & Risk level & PPO order-level ReT & Best comparator & $\Delta$ \\
\midrule
h52 & current & 0.005316 & 0.004832 & +0.000483 \\
h52 & increased & 0.003259 & 0.002517 & +0.000742 \\
h52 & severe & 0.0001193 & 0.0001797 & $-0.0000604$ \\
h104 & current & 0.005648 & 0.005439 & +0.000209 \\
h104 & increased & 0.003660 & 0.003095 & +0.000565 \\
h104 & severe & 0.0001282 & 0.0001188 & +0.0000094 \\
\bottomrule
\end{tabular}
\end{table}

\begin{figure}[htbp]
\centering
\includegraphics[width=0.70\linewidth]{fig5_generalization_heatmap}
\caption{E3 cross-regime/horizon matrix. PPO gains are positive for current and increased risk at both horizons; severe short-horizon disruption is the disclosed boundary case.}
\label{fig:generalization_heatmap}
\end{figure}

The Track B advantage is positive in five of the six evaluated cells,
including the current and increased Garrido-native risk levels at both horizons, with deltas
in several cells (e.g. $h{=}52$, increased: $+0.000742$) exceeding the
canonical \texttt{adaptive\_benchmark\_v2} result reported in
Table~\ref{tab:track_b_main}. The single exception is the shorter horizon
under severe disruption ($h{=}52$), where PPO loses narrowly
($-0.0000604$); this is also the cell with the lowest absolute ReT across
all policies by roughly forty-fold (fill rate $\approx 0.1\%$ for every
evaluated policy), so the comparison sits deep in a service-floor regime
where the difference is not practically meaningful for any policy. At the
canonical horizon under severe risk ($h{=}104$), PPO's win is technically
positive but an order of magnitude smaller than the other four positive
cells and should not be over-interpreted. We report the win as
generalizing broadly rather than universally, with the severe/short-horizon
cell disclosed as a genuine boundary case.

As a further check on whether the disruption regime alone, independent of
the action space, explains the Track B result, we retrain the original
Track A action family (buffer/shift only) directly on
\texttt{adaptive\_benchmark\_v2} --- the same stress regime used for the
Track B canonical result --- at matched scale (5 seeds $\times$ 60{,}000
timesteps). The retrained Track A policy does not outperform the best
constant-fraction static comparator on this regime (Excel ReT 0.00530 vs.\
0.00539; CVaR also favors the static). Access to the stress regime is
therefore not, by itself, sufficient for a learned policy to win; the
result additionally requires action-space coverage of the downstream
bottleneck.

\hypertarget{48-summary}{%
\subsection{Summary}\label{48-summary}}

The experimental evidence forms a coherent narrative across six result
dimensions. Track A demonstrates a structural limitation under
dense-frontier evaluation (Section~4.2). Track B demonstrates that an
operational action-space extension overcomes it under paired
common-random-number evaluation against a 147-cell dense static grid
(Section~4.3). A 5-seed causal ablation localizes the effect to
downstream-dispatch access specifically, ruling out ``more actuators''
as the explanation (Section~4.4). Two independent tests --- an
obs-masked retrain and a privileged-access static comparator --- rule out
that the effect depends on the learned policy exploiting simulator-only
regime information (Section~4.5). Cross-regime and cross-horizon
evaluation shows the advantage generalizes to five of six evaluated
cells, with the one exception (severe disruption at the shorter horizon)
disclosed as a service-floor boundary case, and a companion experiment
shows the disruption regime alone is not sufficient for a win without
downstream-dispatch access (Section~4.7). The central finding is not
that RL is universally effective, but that RL value appears when the
controllable interface reaches the operational bottleneck and the policy
can convert that authority into recovery performance.
